\section{Discussion}
\label{sec:discussion}

In this section, we will discuss selected issue that \greyboxfuzz discover during the evaluation process


\subsection{Unused Variables Issues in Circuit}



Unused public variables may introduce a security risk when a value is declared as a public input and is expected by the verifier to be part of the proved statement, but the circuit does not enforce the intended relationship for that value. Public inputs are visible to the verifier and are commonly used to represent externally checked values, such as an identity commitment, a balance, a nullifier, an output value, or an application-specific parameter. If such a value is declared but not properly connected to the circuit constraints, the verifier may believe that the proof guarantees a statement involving that value, while the circuit proves only a weaker statement. For example, suppose a circuit is intended to prove that a user is older than 18, and the public input is the claimed age. If the circuit declares \texttt{age} as public but only proves knowledge of a private secret and never constrains \texttt{age}, then a malicious prover can generate a valid proof with \texttt{age = 25}, even though the witness does not establish that the user is older than 18. From the verifier's perspective, the proof appears to be associated with the public value \texttt{25}; however, the circuit never checks the age condition. This is not necessarily a failure of proof generation or proof verification. Instead, it is a mismatch between the statement the verifier believes is being proven and the actual constraints enforced by the circuit. The impact can be severe when the unused public variable controls authorization, asset ownership, eligibility, or other security-critical application logic, because an attacker may produce a proof that verifies successfully while bypassing the intended public-input condition.


Currently, \circom can successfully compile a circuit with an unused public variable, and the compiler does not automatically enforce that every public input appears in the intended constraints. This can be dangerous if the proof system does not bind each declared public input to the generated proof, because an attacker may generate a proof for one public input and then reuse the same proof with a different public input. Libraries commonly used with \circom, such as \snarkjs and \rapidsnark, address this input-consistency soundness issue in the R1CS-to-QAP reduction. As described in Remark~2.5 of \cite{ben2014succinct}, some R1CS instances can produce QAP polynomials that are distinct but not linearly independent. The fix increases the QAP degree from \texttt{cs.num\_constraints()+1} to \texttt{cs.num\_constraints()+cs.num\_inputs()+1}, ensuring that public-input polynomials remain independent. Therefore, even if a public variable is unused in the original circuit constraints, the proof is still bound to that declared public input, preventing an attacker from changing the public input after proof generation without invalidating the proof.

Similar fixes have also been adopted by other \zksnark libraries, including \libsnark, \bellman, and \arkworks. However, during our evaluation, \greyboxfuzz found that the unused public variable issue still exists in \playsnark and \gnark \cite{consensys_gnark}. \playsnark is a playground for learning \zksnark proof systems and is not intended for secure deployment. In contrast, \gnark is a production-oriented library that has undergone security audits, but it currently does not follow the R1CS-to-QAP transformation described in Remark~2.5 of \cite{ben2014succinct}. As a result, if an unused public variable exists in the circuit, the generated proof may not be fully bound to that public input. A malicious prover can generate a proof using one public value and then replace that value with another public value during verification, while the verifier may still accept the proof.


\subsection{\zksnark in Federated Learning}
In this work~\cite{han2023kick}, zkSNARKs are used to make the anomaly-detection process in federated learning verifiable. The server acts as the prover and the clients act as verifiers. After the server runs the anomaly-detection mechanism and removes suspected malicious clients, it generates a zkSNARK proof showing that the detection and removal steps were executed correctly. Clients can then verify this proof without redoing the full computation and without relying only on the server's claim. The paper uses zkSNARKs because they provide small proofs and fast verification, which is important when federated-learning clients have limited resources.

The paper uses \zksnark to prove that the server correctly executes the anomaly-detection computation in federated learning. This computation includes several numerical operations, such as matrix multiplication, cosine similarity, division, and square root. Since square root is expensive and difficult to compute directly inside an arithmetic circuit, the paper asks the prover to provide the square-root result \(x\) as a witness and then checks the relationship between \(x\) and the input value \(y\). Instead of computing \(x=\sqrt{y}\) inside the circuit, the circuit checks whether \(x^2 \leq y\) and \((x+1)^2 \geq y\). This design reduces the square-root verification cost to two multiplications and two comparisons.

However, this condition does not exactly prove that \(x\) is the correct integer square root of \(y\). For example, when \(y=9\), the correct square root is \(x=3\), but \(x=2\) also satisfies \(2^2 \leq 9\) and \((2+1)^2 \geq 9\). Therefore, the circuit may accept an incorrect square-root witness for perfect-square values. This is a circuit correctness issue and may become a security issue if the square-root value affects a threshold, distance, standard deviation, anomaly score, or final detection decision. In this paper, the square-root value is used as part of an approximate numerical computation for anomaly detection, so an off-by-one error may have limited impact if the final decision is not sensitive to small numerical changes. However, if a client model is close to the decision boundary, or if the square-root result is used in a threshold-sensitive computation, this approximation may change the anomaly-detection result and should not be considered fully sound.

For this specific application, the \zksnark proof may not require full numerical accuracy because the federated-learning model and anomaly-detection process already tolerate some deviation in model updates and statistical measurements. However, from the cryptographic perspective, the circuit still cannot fully represent the original statement that it intends to prove. Therefore, this issue should be understood as a circuit-level approximation error: it may be acceptable for this particular machine-learning application if it does not affect the final classification result, but it would be unsafe in applications where the square-root value is used in exact security-critical logic.


\section{Future Work}

In general, \greyboxfuzz was designed and developed based on the \zksnark protocol. However, most of its components are not specifically tailored for \zksnark alone and can be extended to support other \zkp protocols. In this section, we discuss the possible extensions of \greyboxfuzz and evaluate its applicability by running several programs on different \zkp protocols.

As presented in \autoref{sec:overview}, \greyboxfuzz consists of four major components: differential fuzzing ground truth, input formatter, coverage monitor, and error locator. The differential fuzzing ground truth and input formatter are designed using Charm, which facilitates cryptographic variable handling in \zkp statements. These components are not tied to a single \zksnark protocol but can represent any cryptographic statement in a \zkp protocol. As a result, no additional implementation or design modifications are required to extend \greyboxfuzz to support differential fuzzing ground truth and input formatting for other \zkp protocols.

The coverage monitor analyzes the binary program and classifies assembly instructions into \cryptobranch and \programbranch using our taint tracking tool. The taint tracking strategy classifies any variable that correlates with input variables as tainted. If this strategy can be applied to another \zkp protocol, no additional implementation is needed for the coverage monitor. For instance, in Bulletproofs, a non-interactive \zkp protocol with short proofs and no trusted setup, the prover takes the circuit, witness, Pedersen commitments, and public parameters (such as $G$ and $H$) as input. The same taint tracking strategy used for \zksnark can be applied to classify \cryptobranch in Bulletproofs since all variables computed based on protocol inputs fall into this category. We believe that most \zkp protocols can adopt a similar strategy to identify \cryptobranch and \programbranch without requiring modifications.

The error locator, which helps identify errors in \zkp programs, requires protocol-specific implementations. A universal error locator is not available for multiple \zkp protocol. However, since the error locator is an optional component of \greyboxfuzz, its absence does not prevent the fuzzer from detecting errors and security vulnerabilities in \zkp binary programs. Thus, even without implementing a dedicated error locator for each \zkp protocol, \greyboxfuzz can still be extended to other \zkp protocols, as all other core components do not require protocol specific modifications.